\documentclass[11pt,letterpaper]{article}

\usepackage[margin=1in]{geometry}
\usepackage{amsmath,amssymb}
\usepackage{booktabs}
\usepackage{listings}
\usepackage{hyperref}
\usepackage{xcolor}
\usepackage[T1]{fontenc}
\usepackage{libertinus-type1}
\usepackage{microtype}

\lstset{
  basicstyle=\ttfamily\small,
  breaklines=true,
  frame=single,
  columns=fullflexible,
  keepspaces=true,
}

\hypersetup{
  colorlinks=true,
  linkcolor=blue!60!black,
  citecolor=blue!60!black,
  urlcolor=blue!60!black,
}

\title{Fixed-Weight Complementary Learning Systems\\on a Ternary Microcontroller:\\The Hippocampus Is the Model}
\author{Tripp Josserand-Austin\\EntroMorphic Research\\\texttt{tripp@entromorphic.com}}
\date{April 12, 2026}

\begin{document}
\maketitle

\begin{abstract}
We present a fixed-weight analog of Complementary Learning Systems (CLS) theory running on a \$0.50 microcontroller. A ternary Closed-form Continuous-time neural network (CfC) with random, non-updating weights serves as the neocortical extractor. A 64-node Navigable Small World graph serves as the hippocampal episodic store. The system classifies 4 wireless signal patterns at 100\% label-free accuracy in peripheral hardware at 430\,Hz, drawing approximately 30 microamps.

The hippocampal layer is permanently load-bearing. We know this not because we chose not to implement consolidation, but because \textbf{we implemented it and it produced no effect.} Three iterations of Hebbian weight learning---with VDB-mismatch targets, TriX-accumulator targets, and diagnosed f-vs-g pathway selection---were tested under ablation-controlled, replicated conditions. Result: $+0.1 \pm 1.1$ Hamming at $n=3$. The neocortex does not learn from the hippocampus.

And yet the hippocampus works. VDB episodic retrieval and CfC feedback blend, with no weight learning, produces 8.5--9.7/80 MTFP divergence across 4 patterns. The VDB is causally necessary: ablation collapses what VDB feedback separates. The transition experiment confirms VDB stabilization label-free: without VDB blend, the old prior reasserts itself; with VDB blend, separation is maintained.

This is the sharpest departure from biological CLS: there is no consolidation, not because we omitted it, but because the system doesn't need it. The hippocampus is not a staging area. It is the temporal model itself.
\end{abstract}

\section{Introduction}

Complementary Learning Systems theory \cite{mcclelland1995} proposes that intelligent systems require two interacting memory systems: a fast learner (hippocampus) that rapidly encodes specific episodes, and a slow learner (neocortex) that gradually extracts statistical structure. The hippocampus enables rapid learning without catastrophic interference in the neocortex. Over time, hippocampal replay trains the neocortex, eventually making the hippocampus redundant for well-consolidated memories.

This paper describes a system that implements the CLS architecture with one critical difference: \textbf{the slow learner never learns---and we tested what happens when it does.}

The CfC weights are random, set at initialization, and never updated during normal operation. There is no replay. There is no consolidation. The hippocampal layer (VDB) is permanently load-bearing.

Unlike prior fixed-weight CLS demonstrations, we do not merely omit consolidation and argue theoretically that the hippocampus is necessary. We implemented the consolidation path---Hebbian LP weight updates from VDB and TriX signals---and measured its contribution under controlled, replicated conditions. The result: no measurable improvement over the hippocampus-only baseline. The zero learning rate is not a design choice. It is the empirically correct operating point.

\subsection{The Architecture}

Three layers, operating on genuine wireless signals from an ESP-NOW sender:

\textbf{Layer 1 (GIE):} A Geometry Intersection Engine running ternary dot products at 430\,Hz on peripheral hardware (GDMA $\to$ PARLIO $\to$ PCNT loopback). Classifies wireless signals at 100\% label-free accuracy (TriX ISR, structural guarantee: $W_f[\text{hidden}] = 0$).

\textbf{Layer 2 (LP core):} A 16\,MHz RISC-V ultra-low-power core ($\sim$30\,$\mu$A) running a 16-neuron ternary CfC with fixed random weights, plus a 64-node NSW vector database storing 48-trit episodic snapshots.

\textbf{Layer 3 (HP core):} The 160\,MHz application core handles initialization, ESP-NOW reception, and test orchestration.

Label-free operation: the pattern\_id field is masked from both TriX enrollment signatures and the GIE runtime input. Classification uses only payload bytes and inter-packet timing features.

\subsection{The CLS Mapping}

\begin{table}[h]
\centering
\begin{tabular}{lll}
\toprule
CLS Component & Biological & The Reflex \\
\midrule
Neocortex & Slow learner & LP CfC: fixed random projection \\
Hippocampus & Fast learner & VDB: 64-node NSW graph \\
Consolidation & Replay trains neocortex & \textbf{Tested. No effect.} $+0.1 \pm 1.1$ ($n=3$) \\
Separation & No interference & $W_f[\text{hidden}] = 0$ \\
Complementarity & Compensates limitations & VDB routes around CfC degeneracies \\
\bottomrule
\end{tabular}
\caption{CLS mapping between biological theory and the Reflex.}
\end{table}

\section{The VDB Is the Temporal Model}

\subsection{The Measurement: MTFP Divergence}

The LP CfC computes 16 integer dot products per step. The MTFP-encoded version (5 trits per dot: sign + 2 exponent + 2 mantissa) is an 80-trit vector that preserves magnitude information. Sign-space collapses the magnitudes, producing LP divergence of 1.2--2.3/16. MTFP-space preserves them, producing 8.5--9.7/80---a richer measurement that correctly identifies effects that sign-space hides.

\subsection{Causal Necessity (TEST 13)}

CMD~4 (CfC step + VDB search, no feedback blend) vs CMD~5 (CfC + VDB + feedback blend):
\begin{itemize}
\item CMD~4 P1-P2 sign Hamming: 1
\item CMD~5 P1-P2 sign Hamming: 2, MTFP Hamming: 9
\end{itemize}

VDB feedback contribution: +1 trit P1-P2 in sign-space. The VDB is causally necessary for LP divergence.

\subsection{The Hippocampal Representation (TEST 12)}

MTFP divergence matrix (/80) after 120 seconds of live wireless input:

\begin{table}[h]
\centering
\begin{tabular}{l rrrr}
\toprule
& P0 & P1 & P2 & P3 \\
\midrule
P0 & 0 & 5 & 6 & 8 \\
P1 & 5 & 0 & 9 & 13 \\
P2 & 6 & 9 & 0 & 10 \\
P3 & 8 & 13 & 10 & 0 \\
\bottomrule
\end{tabular}
\caption{LP MTFP divergence (/80). Mean: 8.5/80. All 6 pairs separated.}
\end{table}

No weight learning. No labels. No training. Just store, retrieve, blend.

\section{Consolidation Tested: Three Iterations of Hebbian Learning}

\subsection{Iteration 1: VDB Mismatch Target, f-Pathway Only}

Per-trit disagreement between VDB best match and current \texttt{lp\_hidden}. Result: $+2.5$ with label in input, $-1.7$ genuinely label-free. The VDB error signal was exploiting pattern\_id leaked through the GIE hidden state.

\subsection{Iteration 2: TriX Accumulator Target, f-Pathway Only}

Clean target (TriX labels are structurally guaranteed). Result: $-1.0$ label-free. $\sim$50\% of errors were in the g-pathway, making f-pathway flips counterproductive.

\subsection{Iteration 3: TriX Accumulator, Diagnosed f+g Pathway}

Per neuron: compare $|f\_\text{dot}|$ vs $|g\_\text{dot}|$, fix the pathway with the smaller dot. Result (3 repetitions, genuinely label-free):

\begin{table}[h]
\centering
\begin{tabular}{l cc}
\toprule
Condition & Sign (/16) & MTFP (/80) \\
\midrule
Control (CMD5 only) & $1.0 \pm 1.3$ & $9.7 \pm 0.6$ \\
Hebbian (CMD5+learn) & $1.1 \pm 0.8$ & $9.7 \pm 1.0$ \\
Contribution & $+0.1 \pm 1.2$ & $+0.1 \pm 1.1$ \\
\bottomrule
\end{tabular}
\caption{Hebbian learning result: indistinguishable from zero at $n=3$.}
\end{table}

\subsection{Why Consolidation Doesn't Help}

The LP CfC has 16 neurons with 48-trit random ternary weights. The Hebbian rule flips one weight per neuron per update, changing the dot product by $\pm 2$. With $\sim$29 non-zero weights per neuron and only a sign-quantized output, many weight configurations produce the same output. The optimization landscape is flat.

The VDB feedback blend directly injects episodic content into \texttt{lp\_hidden}. It bypasses the weights entirely. \textbf{At 16 neurons with random ternary weights, direct episodic injection is more effective than weight learning for producing pattern-discriminative temporal states.}

\section{Kinetic Attention: Harmful at MTFP Resolution}

Agreement-weighted gate bias was tested under label-free conditions. Three runs of Test~14: mean MTFP improvement $-5.5$/80. The bias saturates the GIE hidden state, reducing LP dot magnitude diversity. The sign-space metric ($+1.3$/16) incorrectly showed improvement. Reported as an honest negative result.

\section{The Transition Experiment}

\subsection{Protocol}

TEST~14C: sender transmits P1 for 90 seconds, then switches to P2 for 30 seconds. Three conditions: Full (CMD~5 + gate bias), No bias (CMD~5, bias $= 0$), Ablation (CMD~4, no VDB blend).

Label-free data: Seed~A, \texttt{MASK\_PATTERN\_ID=1 + MASK\_PATTERN\_ID\_INPUT=1}, distinct P2 payload. Multi-seed supporting data: 3 seeds $\times$ 3 conditions (pre-label-free).

\subsection{Ablation Regression: The Hippocampus Stabilizes}

The CLS prediction: removing VDB blend should allow the stale P1 prior to reassert itself post-switch.

\begin{table}[h]
\centering
\begin{tabular}{l rrrrr}
\toprule
Condition & gap@+0 & gap@+5 & gap@+10 & gap@+20 & gap@+30 \\
\midrule
Full & $+12$ & $-2$ & $+18$ & $+14$ & $+14$ \\
No bias & $+14$ & $-3$ & $+10$ & $+11$ & $+2$ \\
Ablation & $+22$ & $+5$ & $-1$ & $+2$ & $\mathbf{-6}$ \\
\bottomrule
\end{tabular}
\caption{MTFP alignment gap (P2 $-$ P1) post-switch, label-free Seed~A. Ablation regresses to $-6$ by step 30.}
\end{table}

The ablation condition regresses: gap starts at $+22$, crosses zero by step~+10, and reaches $-6$ by step~30. The no-bias condition maintains a positive gap. The VDB stabilizes the transition.

TriX accuracy post-switch: 15/15 label-free across all three conditions.

\subsection{Bias Release Dynamics}

\begin{lstlisting}
step +0:  bias=[10, 0]   pred=1  (old prior active)
step +1:  bias=[ 9, 0]   pred=2  (pred flips immediately)
step +5:  bias=[ 5, 0]   pred=2  (geometric decay)
step +10: bias=[ 3, 0]   pred=2
step +20: bias=[ 0,14]   pred=2  (old prior gone, new prior stable)
\end{lstlisting}

The \texttt{pred} flip occurs at step~+1. Old-prior bias decays geometrically ($\times 0.9$/step, half-life $\sim$6.6 steps). New-prior formation begins at step~+15. The prior is a voice that fades, not a verdict that gets revoked.

\section{Related Work}

\subsection{Complementary Learning Systems}

The CLS framework \cite{mcclelland1995} was formalized around the hippocampal-neocortical distinction. Kumaran et al.\ \cite{kumaran2016} updated the theory to account for prefrontal contributions and schema-based learning. O'Reilly and Frank's PBWM model \cite{oreilly2006} extended CLS into the prefrontal-basal ganglia circuit, showing how working memory gating can complement hippocampal episodic encoding---a different form of fast/slow separation, operating at the level of cognitive control rather than memory consolidation.

Hassabis and Maguire \cite{hassabis2007} demonstrated that hippocampal patients cannot construct novel imagined scenes, suggesting the hippocampus is not merely a memory store but a constructive engine. This resonates with our finding: the VDB is not just storing episodes but actively shaping the LP temporal representation through retrieval-and-blend. The temporal model is constructed from memories, not extracted from weights.

Benna and Fusi \cite{benna2016} analyzed memory consolidation in systems with constrained synaptic complexity, showing that consolidation benefits depend on having sufficient degrees of freedom in the slow pathway. Our Hebbian null result ($+0.1 \pm 1.1$ at $n=3$) is consistent with their framework: 16 neurons with $\sim$29 non-zero ternary weights per neuron do not have enough synaptic complexity for consolidation to improve on direct episodic retrieval.

Lengyel and Dayan \cite{lengyel2007} showed that episodic memory systems can outperform statistical learners for decision-making when episodes are few and informative. Our system operates in exactly this regime: 4 patterns, 64 episodes, 120 seconds of accumulation. The VDB-only baseline (9.7/80 MTFP) demonstrates that episodic injection outperforms weight learning in this low-data, low-dimensional setting.

\subsection{Reservoir Computing}

The LP CfC with fixed random ternary weights is functionally a reservoir \cite{jaeger2001echo,maass2002real}. Echo State Networks \cite{jaeger2001echo} and Liquid State Machines \cite{maass2002real} use fixed random recurrent networks as nonlinear projections, training only a linear readout layer. Our system goes further: the readout is also fixed (TriX signatures are enrolled, not trained), and the temporal modeling comes entirely from the episodic memory layer rather than from the reservoir dynamics themselves.

This connects to recent work on untrained random networks as feature extractors. Rosenfeld and Tsotsos \cite{rosenfeld2019intriguing} showed that random-weight CNNs can serve as surprisingly effective feature extractors when combined with appropriate downstream processing. In our system, the ``downstream processing'' is episodic storage and retrieval---not a trained classifier but a memory that accumulates the reservoir's projections over time.

The distinction from standard reservoir computing is that our system has no trained readout. The VDB replaces the readout layer with episodic retrieval, and the temporal representation lives in the memory content rather than in learned output weights. This makes the system fully weight-free in the learning sense: no weights are updated, ever, during operation.

\subsection{Neuromorphic Memory Systems}

On-chip episodic memory for neuromorphic systems has been explored in content-addressable memory architectures \cite{kanerva1988sparse} and associative memory networks \cite{hopfield1982neural}. Kanerva's Sparse Distributed Memory provides a theoretical framework for high-dimensional binary memory systems with similarity-based retrieval. Our VDB NSW graph implements a related function---approximate nearest-neighbor retrieval over ternary vectors---but in 2\,KB of LP SRAM on a commodity microcontroller rather than dedicated associative memory hardware.

\section{The CLS Reframe}

\subsection{The Permanent Hippocampus}

\begin{table}[h]
\centering
\small
\begin{tabular}{lll}
\toprule
CLS Property & Standard Theory & The Reflex \\
\midrule
Hippocampal role & Fast learning, consolidated & \textbf{Permanent temporal model} \\
Neocortical learning & Slow, from replay & \textbf{None. $+0.1 \pm 1.1$ (noise)} \\
Consolidation & Transfer to neocortex & \textbf{Tested. Does not occur.} \\
Redundancy & Eventually redundant & \textbf{Never redundant.} \\
Temporal representation & Neocortical (learned) & \textbf{Hippocampal (episodic)} \\
\bottomrule
\end{tabular}
\caption{Permanent-hippocampus CLS vs standard theory.}
\end{table}

\subsection{When Is the Hippocampus Sufficient?}
\label{sec:sufficiency}

The hippocampus-only path works when three conditions hold:

\begin{enumerate}
\item \textbf{The projection is rough but non-degenerate.} The CfC random projection must separate patterns enough that VDB retrieval returns same-pattern episodes more often than different-pattern episodes. It need not be perfect---the VDB routes around partial degeneracies---but if two patterns project to identical LP trajectories, no amount of episodic content can separate them.

\item \textbf{VDB capacity exceeds the episodic demand.} With 64 nodes and 4 patterns at $\sim$8 inserts per pattern per 90-second run, the VDB fills in $\sim$3 minutes. Each additional pattern requires its own episodic slots. At 8 patterns, the per-pattern allocation drops to $\sim$8 nodes---potentially below the threshold for reliable NSW retrieval ($M=7$ neighbors).

\item \textbf{The task is discriminative, not generative.} The VDB stores and retrieves specific episodes. It does not abstract across them. A task requiring generalization beyond stored states---e.g., classifying a novel pattern not seen during accumulation---would require the weight pathway that our system lacks.
\end{enumerate}

\textbf{Testable predictions for boundary conditions:}

\begin{itemize}
\item \textbf{Pattern scaling:} At 4 patterns (current), the system produces 8.5--9.7/80 MTFP divergence. We predict monotonic degradation as patterns increase, with a critical threshold near $N = 64 / 8 = 8$ patterns where per-pattern VDB allocation drops below the retrieval reliability floor. At $N > 16$, the 48-trit VDB vector may lack the dimensionality to distinguish patterns in NSW Hamming space.

\item \textbf{Projection degeneracy:} Seed~B demonstrates this boundary. Its random weights collapse P1 and P2 to near-parallel LP directions, and VDB retrieval for P1 returns P2 episodes. The system still works (96\% classification from the structural wall) but the temporal model conflates P1 and P2. We predict that projection degeneracy frequency increases with pattern count: more patterns in a 16-dimensional sign space means more collisions.

\item \textbf{Consolidation threshold:} If the LP were widened to 32 or 64 neurons, the Hebbian weight landscape would have more degrees of freedom. There may exist a neuron count above which consolidation begins to improve on episodic retrieval alone. The permanent-hippocampus finding is specific to 16 neurons---it is a capacity claim, not an architectural one.
\end{itemize}

These predictions are currently untested. They define the experimental program that would establish the boundaries of permanent-hippocampus CLS: at what pattern count does it degrade, at what neuron count does consolidation become useful, and what task structure requires generalization beyond episodic retrieval.

\subsection{Implications for CLS Theory}

The ``consolidation makes the hippocampus redundant'' prediction assumes the neocortex can eventually learn what the hippocampus knows. In our system, 16 ternary neurons with $\sim$29 non-zero weights each do not have enough degrees of freedom for Hebbian learning to find useful structure.

This suggests a condition under which consolidation fails: \textbf{when the neocortical substrate has insufficient representational capacity, the hippocampus remains permanently necessary.} This is the normal condition for small embedded systems with fixed-point arithmetic and limited parameters.

\section{Limitations}

\begin{enumerate}
\item \textbf{Transition experiment: single seed label-free.} Multi-seed supporting data confirms ablation regression in Seeds~A and~C (pre-label-free). TriX@15 numbers may differ under label-free conditions.
\item \textbf{Single seed for Hebbian replication.} The $+0.1 \pm 1.1$ finding is from 3 repetitions of one seed. VDB-only baseline ($9.7 \pm 0.6$) is stable across repetitions.
\item \textbf{16 neurons may be too few for consolidation.} Wider LP or MTFP-targeted learning might enable consolidation.
\item \textbf{Four patterns, one sender.}
\item \textbf{JTAG attached.} UART-only verification pending.
\item \textbf{MTFP metric is nonlinear.}
\end{enumerate}

\section{Conclusion}

A ternary microcontroller drawing 30 microamps implements Complementary Learning Systems with a twist: the neocortex never learns---and when we made it try, it couldn't improve on the hippocampus.

Three iterations of Hebbian weight learning produced $+0.1 \pm 1.1$ Hamming improvement over the hippocampus-only baseline. The consolidation path does not work. Not because we omitted it. Because the hippocampus is already doing what consolidation is supposed to achieve.

The VDB episodic memory layer---64 nodes in 2\,KB of LP SRAM, searched at 100\,Hz---stores, retrieves, and blends episodic snapshots. That blend alone produces 8.5--9.7/80 MTFP divergence across 4 wireless signal patterns. The transition experiment confirms: without VDB blend, the old prior reasserts itself; with VDB blend, the system transitions cleanly.

The hippocampus is not a staging area. It is the temporal model itself. The prior should be a voice, not a verdict. And the voice comes from memory, not from learned weights.

\begin{thebibliography}{99}
\bibitem{mcclelland1995}
J.~L.~McClelland, B.~L.~McNaughton, and R.~C.~O'Reilly, ``Why there are complementary learning systems in the hippocampus and neocortex,'' \textit{Psychological Review}, vol.~102, no.~3, pp.~419--457, 1995.

\bibitem{kumaran2016}
D.~Kumaran, D.~Hassabis, and J.~L.~McClelland, ``What learning systems do intelligent agents need? Complementary learning systems theory updated,'' \textit{Trends in Cognitive Sciences}, vol.~20, no.~7, pp.~512--534, 2016.

\bibitem{oreilly2006}
R.~C.~O'Reilly and M.~J.~Frank, ``Making working memory work: A computational model of learning in the prefrontal cortex and basal ganglia,'' \textit{Neural Computation}, vol.~18, no.~2, pp.~283--328, 2006.

\bibitem{hassabis2007}
D.~Hassabis and E.~A.~Maguire, ``Deconstructing episodic memory with construction,'' \textit{Trends in Cognitive Sciences}, vol.~11, no.~7, pp.~299--306, 2007.

\bibitem{benna2016}
M.~K.~Benna and S.~Fusi, ``Computational principles of synaptic memory consolidation,'' \textit{Nature Neuroscience}, vol.~19, no.~12, pp.~1697--1706, 2016.

\bibitem{lengyel2007}
M.~Lengyel and P.~Dayan, ``Hippocampal contributions to control: The third way,'' \textit{Advances in Neural Information Processing Systems}, vol.~20, 2007.

\bibitem{jaeger2001echo}
H.~Jaeger, ``The `echo state' approach to analysing and training recurrent neural networks,'' \textit{GMD Technical Report 148}, German National Research Center for Information Technology, 2001.

\bibitem{maass2002real}
W.~Maass, T.~Natschl\"{a}ger, and H.~Markram, ``Real-time computing without stable states: A new framework for neural computation based on perturbations,'' \textit{Neural Computation}, vol.~14, no.~11, pp.~2531--2560, 2002.

\bibitem{rosenfeld2019intriguing}
A.~Rosenfeld and J.~K.~Tsotsos, ``Intriguing properties of randomly weighted networks: Generalizing while learning next to nothing,'' \textit{Conference on Robot Learning (CoRL)}, 2019.

\bibitem{kanerva1988sparse}
P.~Kanerva, \textit{Sparse Distributed Memory}, MIT Press, 1988.

\bibitem{hopfield1982neural}
J.~J.~Hopfield, ``Neural networks and physical systems with emergent collective computational abilities,'' \textit{Proceedings of the National Academy of Sciences}, vol.~79, no.~8, pp.~2554--2558, 1982.

\bibitem{li2016ternary}
F.~Li, B.~Zhang, and B.~Liu, ``Ternary weight networks,'' \textit{arXiv:1605.04711}, 2016.
\end{thebibliography}

\end{document}
